\section*{Introduction}
This is the report for the SAT model to solve the Street Directionality Problem (SDP).
The idea is to encode the problem with DIMACS formulas and solve it with the \texttt{kissat} SAT solver.
The code is in my public \href{https://github.com/channelEs/CPS_MIRI}{GitHub repository}.
\section*{Variables}

\subsection*{Decision Variables}
\begin{itemize}
    \item \texttt{E\_{i,j}} (
    \textit{Boolean}):
    This variable indicates whether the directed street from crossing $i$ to crossing $j$ is open.
    A value of $1$ means the lane is active, while $0$ means it is closed.
\end{itemize}

\subsection*{Conversion Variables}
\begin{itemize}
    \item \texttt{C\_m} (\textit{Boolean}):
    For each original two-way street $m$, this variable is true when exactly one of its two directions is closed.
    These are the variables used in the objective function.
\end{itemize}

\subsection*{Reachability Variables}
\begin{itemize}
    \item \texttt{R\_{s,i,w}} (\textit{Boolean}):
    This variable represents whether crossing $i$ is reachable from source $s$ within accumulated travel time $w$.
    These variables are used to encode the QoS constraints.
\end{itemize}

\subsection*{Cardinality Variables}
\begin{itemize}
    \item \texttt{S\_{m,r}} (\textit{Boolean}):
    Auxiliary variables used by the sequential counter that enforces the minimum number of conversions.
\end{itemize}

\section*{Constraints}

\subsection*{Street Integrity}
The solver must preserve the original graph structure.
\begin{itemize}
    \item Non-existent streets are forced to $0$.
    \item Original one-way streets are kept fixed.
    \item Original two-way streets must keep at least one direction open.
\end{itemize}

\subsection*{Conversion Encoding}
For each original two-way street, the model links the direction variables with the conversion variable:
\begin{equation}
    C_m \leftrightarrow \neg(E_{u,v} \wedge E_{v,u})
\end{equation}
This means that a conversion is counted exactly when one direction remains open and the other is closed.

\subsection*{Reachability Propagation}
The reachability variables are initialized with the source crossing and then propagated through active streets.
If a node is reachable with some travel time, then it can also reach its neighbors through an open street while respecting the accumulated cost.

\subsection*{Travel Time Limit}
For every source-destination pair $(s,d)$, the final travel time must not exceed the allowed degradation limit:
\begin{equation}
    R_{s,d,\,L_{s,d}} = 1
\end{equation}
where $L_{s,d}$ is the maximum allowed travel time computed from the original shortest path and the percentage $P$.

\subsection*{Cardinality Constraint}
The model uses a sequential counter to enforce a minimum number of converted streets.
The linear search starts from the maximum possible number of conversions and decreases until the formula becomes satisfiable.

\section*{Objective Calculation}
The objective is to maximize the number of converted two-way streets.
If $N$ is the total number of original two-way streets, then the objective can be written as:
\begin{equation}
    \text{maximize } \sum_{m=0}^{N-1} C_m
\end{equation}
Equivalently, the search checks whether it is possible to convert at least $K$ streets, starting from the largest $K$ and going down until SAT is found.

% \section*{Some notes}
% \textbf{Solver strategy:}
% The implementation builds a CNF file, runs \texttt{kissat}, and reads the satisfying assignment from the solver output.

% \textbf{Optimization method:}
% A linear search wrapper is used over the objective value. This makes the SAT approach simple and direct, while still producing an optimal solution when time allows.

% \textbf{Code notes:}
% The SAT implementation is kept separate from the CP and ILP versions, but all three solve the same SDP instance format and produce the same output structure.
